 % !TEX root = ../mainfile.tex
% !TEX encoding = UTF-8 Unicode
% !TEX spellcheck = en_US

\chapter{Literature Review}
\label{ch:literature}

This chapter situates the thesis within six distinct strands of literature. Each body of research addresses a specific dimension of the core question the thesis tries to answer: does the dollar's status as a safe haven depend on the nature of a shock, rather than its source? By reviewing these fields, this chapter establishes the theoretical framework to understand how the dollar functions as a safe haven, setting the stage for examining the consequences of when the United States itself originates a shock, acts as a co-belligerent or is not directly involved.

\section{Currency Risk Factors}
\label{sec:lit_factors}

\textcite{lustig_roussanov_verdelhan_2011} establish the foundational framework for understanding currency risk by identifying two primary common risk factors, a ``level'' factor, the average excess return on foreign-currency portfolios against the dollar, and a ``slope'' factor, which captures the excess returns of carry trades. They demonstrate that high-interest-rate currencies load more heavily on the slope factor, which is closely linked to global equity-market volatility. Crucially, their model shows that during periods of high global risk, low-interest-rate currencies endogenously appreciate, acting as a hedge, while high-interest-rate currencies depreciate.

The dollar's special role is not confined to asset markets. It also dominates the invoicing of international trade. Drawing on evidence from more than 2,500 country pairs covering some ninety per cent of world trade, \textcite{gopinath_boz_casas_diez_gourinchas_plagborgmoller_2020} show that traded-goods prices are predominantly sticky in dollars rather than in the exporter's or the importer's currency. Trade prices and volumes therefore respond to the dollar exchange rate more than to the bilateral rate.

Building on this dominant-currency paradigm, \textcite{dalgic_ozhan_2025} show that a currency's risk premium is not exogenous but rooted in an economy's trade and financial structure, namely the share of its trade invoiced in dollars and the amount of dollar debt it carries. In their small open-economy model, a shock that depreciates the local currency, such as a foreign interest-rate hike, turns contractionary rather than expansionary. Dollar debt opens a financial channel through which the depreciation raises the real value of liabilities. Servicing dollar debt now requires more local currency. Dollar invoicing mutes the trade channel, leaving export volumes largely unresponsive to the exchange rate, providing no offset. Because depreciations are therefore recessionary, the exchange rate turns countercyclical. The currency loses value precisely when domestic income falls, offers its holder no insurance, and is therefore a poor hedge against domestic downturns. Investors hold such currencies only against a sufficient risk premium, the mirror image of the safety commanded by the dollar. The dollar and carry factors, extended to June 2026, form the basis of the portfolio construction in Chapter~\ref{ch:methodology}.

% \textcite{lustig_roussanov_verdelhan_2011} establish the foundational framework for understanding currency risk by identifying two primary common risk factors: a ``level'' factor, representing the dollar's average fluctuations against all foreign currencies, and a ``slope'' factor, which captures the excess returns of carry trades. They demonstrate that high-interest-rate currencies load more heavily on the slope factor, which is closely linked to global equity-market volatility. Crucially, their model shows that during periods of high global risk, low-interest-rate currencies endogenously appreciate, acting as a hedge or safe haven, while high-interest-rate currencies depreciate. This suggests that the dollar's behavior as a safe haven is intrinsically tied to its role as a low-yield currency that provides insurance against global risk by appreciating during adverse shocks.

% That structure is documented by \textcite{gopinath_boz_casas_diez_gourinchas_plagborgmoller_2020}, whose evidence from more than 2,500 country pairs covering some ninety per cent of world trade shows that traded-goods prices are overwhelmingly sticky in dollars rather than in the exporter's or the importer's currency; trade prices and volumes therefore respond to the \emph{dollar} exchange rate rather than to the bilateral rate.

% Building on this dominant-currency paradigm, \textcite{dalgic_ozhan_2025} show that a currency's risk premium is not exogenous but rooted in an economy's trade and financial structure, specifically its share of dollar invoicing and its net foreign (dollar) debt. In their small open-economy model, a shock that depreciates the local currency, such as a foreign interest-rate hike, turns contractionary rather than expansionary. Dollar debt opens a financial channel through which the depreciation raises the real value of liabilities, while dollar invoicing mutes the trade channel, leaving export volumes largely unresponsive to the exchange rate. Because depreciations are therefore recessionary, the exchange rate becomes countercyclical, meaning the currency depreciates during domestic reseccions, losing value when income is lowest, making it a poor hedge, and hence a risky currency to hold. Investors are compensated for holding such currencies with a standing risk premium. For the empirical design, this strand supplies the building blocks: the dollar and carry factors, rebuilt through mid-2026, underpin the portfolio construction of Chapter~\ref{ch:methodology}.

\section{Dollar Dominance and Its Erosion}
\label{sec:lit_dominance}

The dollar's role as the primary global reserve currency confers an ``exorbitant privilege'', a persistent excess return on U.S.\ external assets over its liabilities \parencite{gourinchas_rey_2005}. \textcite{gourinchas_rey_govillot_2017} show that this privilege is mirrored by an ``exorbitant duty''. In global crises the United States transfers wealth abroad, roughly $19\%$ of U.S.\ GDP during 2007--2009. Because it holds risky foreign assets while owing safe dollar liabilities, collapsing global asset prices and the appreciating dollar both push its net foreign asset position downward, with the exchange rate accounting for roughly a third of the valuation effect. \textcite{jiang_krishnamurthy_lustig_2021} find the source of that movement, showing that surges in foreign demand for safe dollar assets appreciate the dollar immediately. They capture the demand through the convenience yield on U.S.\ Treasuries, the return investors give up in exchange for safety and liquidity. The dollar's appreciation in a global crisis is therefore a consequence of its role as the world's supplier of safe assets. \textcite{obstfeld_zhou_2023} add the cyclical dimension, describing a ``global dollar cycle'' in which dollar appreciation tightens financial conditions worldwide. They find that dollar movements correlate most reliably with measures of investor risk appetite.

While the 2008 financial crisis and the COVID-19 pandemic elicited the standard appreciation, \textcite{jiang_krishnamurthy_lustig_richmond_2026} argue that the dollar has been losing that standing gradually, documenting convenience yields on dollar safe assets eroding over the two years before the 2025 tariff announcements. A comparable trend is evident in central bank reserves. \textcite{arslanalp_eichengreen_simpsonbell_2022} document a two-decade decline in the dollar's share of global reserves, and attribute it to deliberate diversification by reserve managers rather than to valuation effects. Between 1999 and 2021 the dollar's share of global reserves declined from 71\% to 59\%. A quarter of what left the dollar moved into the Chinese renminbi, the remaining three quarters into currencies of smaller economies that had rarely served as reserve assets before.

% The dollar's role as the primary global reserve currency confers an ``exorbitant privilege''---a persistent excess return on U.S.\ external assets over its liabilities \parencite{gourinchas_rey_2005}---mirrored by an ``exorbitant duty'': in global crises the United States transfers wealth to the rest of the world, roughly $19\%$ of U.S.\ GDP during 2007--2009, as the dollar's flight-to-safety appreciation revalues its external balance sheet \parencite{gourinchas_rey_govillot_2017}. On the mechanism, \textcite{jiang_krishnamurthy_lustig_2021} show that surges in foreign demand for safe dollar assets, measured through the convenience yield embedded in the Treasury basis, appreciate the dollar immediately. This is precisely the reflex a safe-haven currency should display.

% Crucially, this literature implies that the haven response depends on what a shock does to the demand for dollar safety. While the 2008 financial crisis and the COVID-19 pandemic elicited the standard appreciation, \textcite{jiang_krishnamurthy_lustig_richmond_2026} document an anomalous depreciation after the 2025 U.S.\ tariff announcements---a ``flight away'' from the dollar examined in detail in Section~\ref{sec:lit_tariffs}. \textcite{obstfeld_zhou_2023} add the cyclical dimension: a ``global dollar cycle'' in which dollar appreciation tightens financial conditions worldwide, with dollar movements correlating most reliably with investor risk appetite. A slow-moving counterpart to any abrupt break is visible in official reserves, which \textcite{arslanalp_eichengreen_simpsonbell_2022} show have been actively diversified away from the dollar for two decades. For the thesis, this strand establishes the stakes: if the dollar's haven status is structural, it should hold across shock types; if it is eroding, U.S.-originated policy shocks are where the cracks should show first.

\section{Safe-Haven Currencies and Gold}
\label{sec:lit_safe_haven}

A central paradox, identified by \textcite{habib_stracca_2012}, is the dollar's appreciation during the 2008 global financial crisis, a crisis that originated in the United States. Analysing 52 currencies over almost 25 years, they find that safe-haven status is best predicted not by the interest-rate spread that the carry-trade literature emphasises but by a country's net foreign asset position, which signals how exposed it is to external stress. A currency that has hedged well in the past also tends to continue doing so, and the currencies of large but less financially open economies hedge better in periods of distress. The interest-rate spread against the United States, by contrast, matters only during the 2008 crisis itself.

\textcite{cheema_faff_szulczyk_2022} compare a range of safe-haven assets across different crises, and find their safe-haven character changed considerably during the COVID-19 pandemic compared with the 2008 financial crisis. The Swiss franc remained an intermediate haven in most settings, while gold and silver, weak havens in 2008, turned risky. The dollar deteriorated as well, from intermediate protection in 2008 to weak protection during the pandemic.

Providing broader context, \textcite{debock_decarvalhofilho_2015} demonstrate that the dollar, Japanese yen and Swiss franc consistently appreciate against other currencies during risk-off episodes. How much each currency gains depends on its interest rate before the episode, its net foreign asset position, whether it is overvalued, and how cheaply and freely it can be traded. \textcite{ranaldo_soderlind_2010} show that during U.S.\ equity sell-offs the franc and yen appreciate against the dollar, with effects that appear within hours, persist for several days, and grow more than proportionally as volatility rises. The dollar therefore strengthens broadly in risk-off episodes while still losing ground to the franc and the yen.

% A similar conditionality holds for gold. \textcite{baur_lucey_2010} find that gold hedges equities on average, but that it works as a safe haven only after severe market declines, and even then only for a short time. This strand supplies the thesis's working definition of a safe haven and its benchmarks. Chapter~\ref{ch:framework} adopts the conditional, crisis-contingent notion of safety, and the yen, the franc and gold serve as the standards against which the dollar's response is measured.

\textcite{baur_lucey_2010} are the first to formally define and test the distinction between a hedge and a safe haven. In their terminology, a hedge is uncorrelated or negatively correlated with equities on average, which does not guarantee protection during market crashes. A safe haven must deliver exactly that protection in the extreme. Applying the distinction to gold, they find that it hedges equities on average, but works as a safe haven only after severe market downturns, and even then only for a short time. This strand supplies the thesis's working definition of a safe haven and its benchmarks. Chapter~\ref{ch:framework} adopts the conditional, crisis-contingent notion of safety, and the yen, the franc and gold serve as the standards against which the dollar's response is measured.

% The literature on safe-haven currencies indicates that the U.S.\ dollar's status is not static but varies with the nature and epicentre of a shock. A central paradox, identified by \textcite{habib_stracca_2012}, is the dollar's appreciation during the 2008 global financial crisis---an episode the United States itself exported. They find that while the net foreign asset position is the most robust long-term predictor of safe-haven status, the dollar's 2008 behaviour was anomalous compared with historical regularities.

% The comparative efficacy of the dollar across shocks is most directly addressed by \textcite{cheema_faff_szulczyk_2022}, who find that the dollar's safe-haven character weakened considerably during the COVID-19 pandemic compared with 2008: an intermediate haven in the earlier crisis, it performed weakly across most settings in the later one---the empirical counterpart of the demand-side reading of the same episodes in Section~\ref{sec:lit_dominance}. Providing broader context, \textcite{debock_decarvalhofilho_2015} demonstrate that the dollar, Japanese yen and Swiss franc consistently appreciate during risk-off episodes, with magnitudes shaped by pre-crisis yields and market-liquidity factors. At the highest frequencies, however, the ordering within pairs reverses: \textcite{ranaldo_soderlind_2010} show that in U.S.\ equity sell-offs the franc and yen appreciate \emph{against the dollar}, with these haven properties materialising within hours to days and strengthening non-linearly in volatility---so against the classic havens themselves, the dollar is the risky leg. Gold supplies the natural non-sovereign benchmark: \textcite{baur_lucey_2010} find that gold is a hedge against equities on average but a safe haven only in extreme declines, and only briefly. This conditional, crisis-contingent notion of safety is the one the thesis adopts (Chapter~\ref{ch:framework}) and tests for the dollar, with the yen, the franc and gold as the benchmarks against which its response is judged.

\section{Geopolitical Risk}
\label{sec:lit_gpr}

% Recent research suggests that the dollar's response to geopolitical shocks is contingent on the nature of the event and a country's position in the global network. \textcite{caldara_iacoviello_2022} provide the foundational measurement, distinguishing geopolitical ``threats''---the expectation of future adverse events---from geopolitical ``acts'', their actual realisation; a shock to their news-based index induces persistent declines in investment, employment and stock prices, driven by both the threat and the realisation components. Expanding this to exchange rates, \textcite{yilmazkuday_2025} demonstrates that the dollar's geopolitical safe-haven status is not universal: shocks to global geopolitical risk---driven primarily by geopolitical acts---produce a statistically significant depreciation of the U.S.\ dollar, even as currencies such as the South African rand or Brazilian real appreciate, with the currencies of countries more deeply embedded in global value chains depreciating most. Complementing this, \textcite{liu_zhang_2024} find that a strategy long high-GPR currencies and short low-GPR currencies earns a significant excess return, and that this return predictability stems from country-specific and regional risk components rather than a single global factor---implying that a currency's risk--return profile is sensitive to the geographic origin of the shock.

% Institutional evidence points in the same direction: the \textcite{imf_2025_gfsr} finds that although geopolitical risk events exert only a modest effect on asset prices on average, major events---military conflicts in particular---produce substantially larger and more persistent equity declines and raise sovereign risk premia, with the impact concentrated in economies holding weaker fiscal and external buffers and transmitted across borders through trade and financial linkages.

% The thesis draws on this strand directly: the \textcite{caldara_iacoviello_2022} index dates and scales the Hormuz and Ukraine episodes, its threat--act distinction maps onto the staged escalation of the Hormuz crisis, and the GPR series serves as an external control in the empirical chapters.

\textcite{caldara_iacoviello_2022} provide the foundational measurement of
geopolitical risk, a news-based index that separates geopolitical ``threats'',
events that are feared but have not yet occurred, from geopolitical ``acts'', their
actual realisation. Shocks to the index depress investment, employment and stock
prices, and the damage persists, with both the threat and the realisation
components contributing. Expanding this to exchange rates,
\textcite{yilmazkuday_2025} demonstrates that the dollar's geopolitical safe-haven
status is not universal. A year after a shock to global geopolitical risk, the U.S.\ dollar has depreciated significantly, even as currencies such as the South African rand or Brazilian real have appreciated. The currencies of countries most deeply
integrated into international supply chains depreciate most, a pattern driven
primarily by geopolitical acts rather than threats. Complementing this,
\textcite{liu_zhang_2024} find that a zero-cost strategy that buys the currencies
of high-geopolitical-risk countries and sells those of low-risk countries earns a
significant excess return of 5.72\% per year, and that this return predictability
stems from country-specific and regional risk components rather than a single
global factor. A currency's risk--return profile is therefore sensitive to the
geographic origin of the shock.

Institutional evidence points in the same direction. The \textcite{imf_2025_gfsr}
finds that although geopolitical risk events exert only a modest effect on asset
prices on average, major events, military conflicts in particular, produce
substantially larger and more persistent equity declines and raise governments'
borrowing costs. The effects are strongest in emerging markets with little fiscal
room and limited foreign reserves, and they spill over to other countries through
trade and financial ties.

The thesis draws on this strand directly. The index places the Hormuz and Ukraine
episodes in their longer-run context in Chapter~\ref{ch:background}, its threat and
act components are tracked around each shock in the event study, and changes in the
index serve as a measure of geopolitical risk in the safe-haven regression.

\section{Tariff Shocks and Exchange Rates}
\label{sec:lit_tariffs}

% The 2025 tariff episode has generated a fast-growing literature that is central to this thesis. \textcite{kamin_2025} provides econometric evidence that, following the Liberation Day announcements of April~2025, the dollar switched from behaving as a safe haven---appreciating in periods of market volatility---to behaving as a ``risk-on'' currency that moves inversely with volatility. \textcite{jiang_krishnamurthy_lustig_richmond_2026} document the same break in historical correlations---the dollar depreciating while U.S.\ rates rose relative to the euro, the VIX increased, and the dollar convenience yield fell---and interpret it through shifts in global demand for dollar safe assets, calibrating a steady-state real dollar depreciation of around $7.6\%$ should the United States lose reserve-currency demand. \textcite{hassan_mertens_wang_zhang_2025} provide the theoretical counterpart: in a general-equilibrium model in which currency safety is endogenous, a trade war that isolates U.S.\ goods markets erodes the dollar's safety premium and raises U.S.\ interest rates.

% On the empir ical exchange-rate response, \textcite{ostry_lloyd_corsetti_2025} construct a tariff-shock database and show that exchange rates react to U.S.\ tariff shocks systematically once the size and relevance of the measures---and expected retaliation---are accounted for. Notably, the 2025 depreciation inverts the pre-2025 evidence, in which U.S.\ tariff actions tended to \emph{strengthen} the dollar: \textcite{khalil_strobel_2024} trace the dollar's appreciation during the 2018--2019 tariffs to trade policy itself, making the 2025 episode a genuine puzzle.

% \textcite{kalemliozcan_soylu_yildirim_2026} resolve that puzzle by showing the standard prediction is overturned once policy \emph{uncertainty} is taken seriously: in a two-country general-equilibrium model with risk-averse agents and segmented financial markets, tariff volatility enters uncovered interest parity as a risk-premium wedge, so that higher tariff uncertainty raises precautionary saving and risk premia and depreciates the currency immediately, even as tariffs themselves rise---a mechanism they calibrate to reproduce the size and timing of the observed 2025 dollar depreciation.

% Closest to the present design, \textcite{benguria_saffie_2025} run an event study of the 2~April 2025 announcement and identify sizeable declines in U.S.\ equity prices that scale with the tariff increment. This strand supplies the thesis's first treatment: the Liberation Day shock, whose event-study comparison against contemporaneous military/oil-supply shocks is the contribution of this thesis.

The 2025 tariff episode has generated a fast-growing literature that is central to
this thesis. \textcite{kamin_2025} estimates the dollar's daily relationship with the VIX and interest differentials and finds a significant break after the Liberation Day announcements of
April~2025. Its sensitivity to the VIX turns from positive to negative, so a currency that
once gained when volatility spiked now falls with it, and investors no longer
treated the dollar as a refuge in stress.
\textcite{jiang_krishnamurthy_lustig_richmond_2026} document the same break across
a wider set of correlations. The dollar depreciated while U.S.\ interest rates rose
relative to euro rates, the VIX climbed and the convenience yield on Treasuries
fell, a pattern they attribute to weakening foreign demand for the safe assets the
United States supplies. \textcite{hassan_mertens_wang_zhang_2025} provide the
theoretical counterpart in a general-equilibrium model where a currency's safety
arises endogenously. In their account the dollar is safe because disturbances originating in the United States move traded-goods prices worldwide, and a trade war that cuts the country off from world markets shuts down those spillovers, so the dollar's habit of strengthening in bad times fades and U.S.\ interest rates move higher.

The depreciation was not what the earlier record suggested.
\textcite{khalil_strobel_2024} show that the tariff rounds of 2018--2019 came with
a stronger dollar and trace that appreciation to U.S.\ trade policy itself.
\textcite{ostry_lloyd_corsetti_2025} qualify the pattern with a database of tariff
announcements, threats and implementations spanning 2018--2020 and 2025, weighting
each shock by the size of the tariff and its economic relevance. The direction of
the response, they find, hinges on retaliation. The dollar appreciates when the
United States imposes tariffs unilaterally but depreciates when trading partners
strike back, which leads them to conclude that the fall after 2~April 2025 was no
anomaly once retaliation is taken into account.
\textcite{kalemliozcan_soylu_yildirim_2026} offer a complementary explanation
centred on uncertainty. In their model, doubt about future tariff policy feeds a
risk premium into the exchange rate, and greater tariff uncertainty raises
precautionary saving and the compensation investors demand, weakening the currency
at once even though the tariffs alone would strengthen it. Calibrated to the 2025
episode, the model matches how far and how fast the dollar fell.

\textcite{benguria_saffie_2025} measure the immediate market response to the
2~April announcement itself, separating the anticipated part of each country's
tariff from the surprise created by rounding the rates upward. Only the surprise
moves prices, a one-percentage-point higher tariff lowering U.S.\ equity prices by
0.23\%, and the dollar shows no appreciation, if anything a depreciation against
floating currencies. This strand supplies the thesis's first treatment. The
Liberation Day shock enters the event study as the trade-policy episode against
which the military and oil-supply shocks are compared.

\section{Commodity Currencies and Oil}
\label{sec:lit_commodities}

% \textcite{chen_rogoff_2003} establish that the U.S.-dollar price of commodity exports exerts a strong and stable influence on the real exchange rates of commodity exporters such as Australia, Canada and New Zealand, providing the basis for the oil-exporter versus oil-importer portfolio split used in Chapter~\ref{ch:methodology}. The nature of the oil shock matters, however: \textcite{kilian_2009} shows that oil-price movements driven by supply disruptions, aggregate demand, and precautionary (oil-specific) demand have distinct macroeconomic consequences, so a price increase caused by a supply scare is not equivalent to one caused by strong demand---nor to a tariff-induced demand collapse. This typology maps directly onto the thesis's contrast between the demand-side tariff shock (falling oil) and the supply-side Hormuz shock (surging oil). The oil--dollar link also runs through the exchange rate directly: \textcite{chen_liu_wang_zhu_2016} show that the dollar's response to oil-price shocks depends on whether they originate in supply or demand. The anatomy of the 2026 Hormuz episode itself, and the sources documenting it, are set out in Chapter~\ref{ch:background}.

% \textcite{chen_rogoff_2003} establish the commodity-currency link for Australia, Canada and New Zealand. The world price of their commodity exports, which these small economies take as given, feeds directly into their real exchange rates, most clearly for Australia and New Zealand. This link provides the basis for the oil-exporter versus oil-importer portfolio split used in Chapter~\ref{ch:methodology}. The nature of the oil shock matters, however. \textcite{kilian_2009} decomposes oil-price movements into supply disruptions, shocks to global demand for industrial commodities, and precautionary demand driven by concern about future supply, and shows that each moves the price of oil and the macroeconomy differently. A price increase caused by a supply scare is therefore not equivalent to one caused by strong demand, nor to a tariff-induced demand collapse. This typology maps directly onto the thesis's contrast between the demand-side tariff shock, with falling oil, and the supply-side Hormuz shock, with surging oil. The oil--dollar link also runs through the exchange rate directly. \textcite{chen_liu_wang_zhu_2016} show that the dollar's response to oil-price shocks depends on whether the shocks originate in supply or demand. The anatomy of the 2026 Hormuz episode itself, and the sources documenting it, are set out in Chapter~\ref{ch:background}.

\textcite{chen_rogoff_2003} show that for commodity exporters such as Australia, Canada and New Zealand, the world price of their exports is a strong driver of the real exchange rate. These three serve as the natural test cases, developed economies with floating exchange rates whose exports remain dominated by commodities. The effect is clearest for the Australian and New Zealand dollars. \textcite{kilian_2009} decomposes oil-price changes into supply disruptions, global demand for industrial commodities, and precautionary demand reflecting concern about future supply, each with distinct effects on the oil price, output and inflation. A supply-driven price increase is therefore not equivalent to a demand-driven one.

\textcite{chen_liu_wang_zhu_2016} show that the dollar reacts more strongly to demand-driven than to supply-driven oil shocks. The reaction is largest against the Canadian dollar and the Norwegian krone, because Canada and Norway trade oil heavily with the United States. Oil has also become more important for the dollar over time. Before the global financial crisis, oil shocks accounted for 10\% to 20\% of the movement in the dollar's bilateral rates. Afterwards the share rises to over 40\%, and for the Canadian dollar to more than half.

\textcite{albulescu_ajmi_2021} document predictive causality between oil prices and the dollar in both directions, strongest in crisis periods. \textcite{salisu_olaniran_tchankam_2022} show that oil tail risk predicts the tail risk of the dollar's rates against the Canadian dollar and the British pound, while the dollar-yen pair moves inversely, consistent with the yen's haven role.

For the thesis, this literature provides the basis for the oil-exporter versus oil-importer portfolio split in Chapter~\ref{ch:methodology}, and the distinction between supply-driven and demand-driven oil movements \parencite{kilian_2009} separates the Hormuz crisis from the tariff shock. The episodes themselves, and the sources that document them, are covered in Chapter~\ref{ch:background}.

% \section{Synthesis and Research Gap}
% \label{sec:lit_synthesis}

% Across these six strands, the dollar's safe-haven status has been established structurally, as a feature of the international monetary system; conditionally, as a crisis-contingent property in the sense of \textcite{baur_lucey_2010}; and, most recently, as something the 2025 trade-policy shock may have disturbed. What the literature has not yet done is hold a U.S.-originated trade-policy shock against a contemporaneous military and oil-supply shock within a single research design, carrying the oil channel and a gold benchmark through the same lens. That comparison is the gap this thesis addresses. The next chapter distils these strands into a working definition of a safe-haven currency, a classification that separates a shock's origin from its nature, and the hypotheses that follow.

% Across these six strands, the dollar's safe-haven status appears in three guises. It is structural, a feature of the international monetary system. It is conditional, a crisis-contingent property in the sense of \textcite{baur_lucey_2010}. And it is contested, since the 2025 trade-policy shock may have disturbed it. The remaining strands supply the tools for examining it, the factor and portfolio framework, the geopolitical-risk index, and the oil channel that separates supply-driven from demand-driven shocks. What the literature has not yet done is hold a U.S.-originated trade-policy shock against a contemporaneous military and oil-supply shock within a single research design, carrying the oil channel and a gold benchmark through the same analysis. That comparison is the gap this thesis addresses. The next chapter distils these strands into a working definition of a safe-haven currency, a classification that separates a shock's origin from its nature, and the hypotheses that follow.